% Part: computability
% Chapter: recursive-functions
% Section: primitive-recursion

\documentclass[../../../include/open-logic-section]{subfiles}

\begin{document}

\olfileid{cmp}{rec}{pre}
\olsection{ఆదిమ పునరావృత్తి}

ఉత్తరగామి క్రియ~$+1$ను పరిమిత సంఖ్యలో ప్రయోగించడం ద్వారా ప్రతి సహజ
సంఖ్యనూ~$0$ నుంచి చేరవచ్చనేది సహజ సంఖ్యల ఒక లక్షణం---ఏ సహజ సంఖ్యైనా
~$0$ లేదా …~$0$ యొక్క ఉత్తరగామికి ఉత్తరగామి. ఈ వాస్తవాన్ని
ఉపయోగించుకొనే $h\colon\Nat \to \Nat$ ప్రమేయాన్ని నిర్దేశించే ఒక
పద్ధతి ఇది: (a)~$0$ ఆర్గుమెంటుకు $h$ విలువ ఏమిటో నిర్దేశించాలి;
(b)~$h(x)$ విలువ ఇచ్చినప్పుడు $h(x+1)$ విలువను ఎలా గణించాలో కూడా
నిర్దేశించాలి. (a) మనకు $h(0)$ ఏమిటో నేరుగా చెబుతుంది కనుక $h$,
~$0$ వద్ద నిర్వచితమవుతుంది. ఇప్పుడు $x=0$కు (b) ఇచ్చిన సూచనను
ఉపయోగించి, $h(0)$ నుంచి $h(1)=h(0+1)$ను గణించవచ్చు. అదే సూచనను
$x=1$కు ఉపయోగించి, $h(1)$ నుంచి $h(2)=h(1+1)$ను గణిస్తాం; ఇదే విధంగా
కొనసాగిస్తాం. ప్రతి సహజ సంఖ్య~$x$కూ, $h(x)$ నుంచి $h(x+1)$ను
నిర్వచించే దశను చివరకు చేరుతాం; అందువల్ల ప్రతి $x\in\Nat$కూ $h(x)$
నిర్వచితమవుతుంది.
\sourcecorrection{OLTECMPREC-001}{ఆంగ్ల మూలం చివరి వాక్యంలో దిశను తారుమారు చేసి $h(x+1)$ నుంచి $h(x)$ను నిర్వచిస్తామని చెప్పింది; ముందరి వివరణకు అనుగుణంగా $h(x)$ నుంచి $h(x+1)$ను నిర్వచించేలా సరిచేశాం.}

ఉదాహరణకు, కింది రెండు సమీకరణాలతో $h\colon \Nat \to \Nat$ను
నిర్దేశించామనుకుందాం:
\begin{align*}
h(0) & =  1\\
h(x+1) & =  2 \cdot h(x)
\end{align*}
గుణించడం ఎలాగో మనకు ఇప్పటికే తెలిస్తే, పై (a), (b)లకు కావలసిన
సమాచారాన్ని ఈ సమీకరణాలు ఇస్తాయి. రెండవ సమీకరణాన్ని వరుసగా
ప్రయోగించడం ద్వారా మనకు
\begin{align*}
  h(1) & = 2\cdot h(0) = 2,\\
  h(2) & = 2\cdot h(1) = 2\cdot 2,\\
  h(3) & = 2 \cdot h(2) = 2\cdot 2 \cdot 2,\\
  & \vdots
\end{align*}
వస్తాయి. మనం నిర్దేశించిన ప్రమేయం~$h$కు $h(x)=2^x$ అని తెలుస్తుంది.

ఈ రెండు ప్రమాణాలను తీర్చే ప్రమేయం~$h$ ఒక్కటే ఉంటుందని సహజ సంఖ్యల
లక్షణం హామీ ఇస్తుంది. ఇలాంటి సమీకరణాలు $h$ను నిర్వచిస్తాయి; వాటి జంటను
ప్రమేయం~$h$కు
\emph{ఆదిమ పునరావృత్తి ద్వారా నిర్వచనం} అంటారు. దీన్ని
``పునరావృత్త'' నిర్వచనం అనడానికి కారణం, నిర్వచనంలో---ప్రత్యేకించి
రెండవ సమీకరణం కుడివైపున---$h$ స్వయంగా రావడం. $h(x+1)$ను
నిర్వచించేటప్పుడు వెంటనే ముందున్న $h(x)$ విలువను మాత్రమే వాడతాం కనుక
ఇది ``ఆదిమం.'' $\Nat$పై ఒక ప్రమేయాన్ని పునరావృత్తంగా నిర్వచించే
అత్యంత సరళ పద్ధతి ఇదే.

కూడిక, గుణకారం వంటి ఇంకా ప్రాథమికమైన ప్రమేయాలను కూడా ఆదిమ
పునరావృత్తి ద్వారా నిర్వచించవచ్చు. అయితే ఈ సందర్భాల్లో సంబంధిత
ప్రమేయాలు $2$-స్థానికాలు. ఒక ఆర్గుమెంటు స్థానాన్ని స్థిరపరచి, మరొక
స్థానాన్ని పునరావృత్తికి వాడతాం. ఉదాహరణకు $\Add(x,y)$ను
నిర్వచించడానికి, $x$ను స్థిరపరచి మొదట $y=0$కు విలువను, తరువాత $y$
విలువ ఆధారంగా $y+1$కు విలువను నిర్వచించవచ్చు. $x$ స్థిరంగా ఉంది కనుక
నిర్వచక సమీకరణాల ఎడమ, కుడి వైపుల రెండింటిలోనూ కనిపిస్తుంది.
\begin{align*}
\Add(x,0) & =  x\\
\Add(x,y+1) & =  \Add(x,y)+1
\end{align*}
ఈ సమీకరణాలు అన్ని $x$, $y$లకు $\Add$ విలువను నిర్దేశిస్తాయి.
ఉదాహరణకు $\Add(2,3)$ను కనుగొనడానికి, $x=2$గా నిర్వచక సమీకరణాలను
ప్రయోగిస్తాం: మొదటిదానితో $\Add(2,0)=2$ను, తరువాత రెండవదానితో వరుసగా
$\Add(2,1)=2+1=3$, $\Add(2,2)=3+1=4$, $\Add(2,3)=4+1=5$ను
కనుగొంటాం.

$\Add$ నిర్వచనంలోని రెండవ సమీకరణం కుడివైపున $+$ను వాడాం; కానీ
~$1$ను కూడడానికి మాత్రమే వాడాం. మరో మాటలో, ఉత్తరగామి ప్రమేయం
$\Succ(z)=z+1$ను వాడి, ముందరి విలువ $\Add(x,y)$కు దాన్ని ప్రయోగించి
$\Add(x,y+1)$ను నిర్వచించాం. కనుక, ముందరి విలువకు ప్రయోగించే ఒకే
ప్రమేయం పరంగా ఈ పునరావృత్త నిర్వచనాన్ని ఆలోచించవచ్చు. అయితే ఆ
ప్రమేయం ముందరి విలువపైనే కాక $x$, $y$లపైనా ఆధారపడేందుకు అనుమతించడం
నష్టం చేయదు---కొన్నిసార్లు అది అవసరమే. పరిశీలించండి:
\begin{align*}
  \Mult(x,0) & =  0 \\
  \Mult(x,y+1) & =  \Add(\Mult(x,y),x)
\end{align*}
ఇది ముందరి విలువ $\Mult(x,y)$కూ మొదటి ఆర్గుమెంటు~$x$కూ $\Add$
ప్రమేయాన్ని ప్రయోగించడం ద్వారా ఇచ్చిన $\Mult$ యొక్క ఆదిమ పునరావృత్త
నిర్వచనం. అన్ని $x$, $y$ ఆర్గుమెంటులకు $\Mult(x,y)$ను కూడా ఇది
నిర్వచిస్తుంది. ఉదాహరణకు, $\Mult(2,0)$, $\Mult(2,1)$,
$\Mult(2,2)$, $\Mult(2,3)$లను వరుసగా గణించడం ద్వారా $\Mult(2,3)$
నిర్ణయమవుతుంది:
\begin{align*}
  \Mult(2,0) & = 0\\
  \Mult(2,1) & = \Mult(2,0+1) =
  \Add(\Mult(2,0), 2) = \Add(0, 2) = 2\\
  \Mult(2,2) & = \Mult(2,1+1) =
  \Add(\Mult(2,1), 2) = \Add(2, 2) = 4\\
  \Mult(2,3) & = \Mult(2,2+1) =
  \Add(\Mult(2,2), 2) = \Add(4, 2) = 6
\end{align*}

సాధారణ నమూనా ఇది: $h(x_0,\dots,x_{k-1},y)$ ప్రమేయానికి ఆదిమ
పునరావృత్త నిర్వచనం ఇవ్వడానికి రెండు సమీకరణాలను ఇస్తాం. మొదటిది
$h$ను ప్రస్తావించకుండా $h(x_0,\dots,x_{k-1},0)$ విలువను
నిర్వచిస్తుంది. రెండవది $h(x_0,\dots,x_{k-1},y)$, ఇతర ఆర్గుమెంటులు
$x_0$, \dots,~$x_{k-1}$, అలాగే~$y$ పరంగా
$h(x_0,\dots,x_{k-1},y+1)$ విలువను నిర్వచిస్తుంది. ఆ రెండవ
సమీకరణంలో $h$ యొక్క వెంటనే ముందరి విలువను మాత్రమే వాడవచ్చు. ఈ రెండు
సమీకరణాల కుడివైపులు ఇచ్చే క్రియలను $f$, $g$ అనే ప్రమేయాలుగా
భావిస్తే, ఆదిమ పునరావృత్తి ద్వారా కొత్త ప్రమేయం~$h$ను నిర్వచించే
సాధారణ నమూనా ఇది:
\begin{align*}
  h(x_0, \dots, x_{k-1}, 0) & = f(x_0, \dots, x_{k-1})\\
  h(x_0, \dots, x_{k-1}, y+1) & =
  g(x_0, \dots, x_{k-1}, y, h(x_0, \dots, x_{k-1}, y))
\end{align*}
$\Add$ సందర్భంలో $k=1$, $f(x_0)=x_0$ (తాదాత్మ్య ప్రమేయం),
$g(x_0,y,z)=z+1$ (తన మూడవ ఆర్గుమెంటు ఉత్తరగామిని ఇచ్చే $3$-స్థానిక
ప్రమేయం):
\begin{align*}
  \Add(x_0, 0) & = f(x_0) = x_0\\
  \Add(x_0, y+1) & = g(x_0, y, \Add(x_0, y)) =
  \Succ(\Add(x_0, y))
\end{align*}
$\Mult$ సందర్భంలో $f(x_0)=0$ (ఎల్లప్పుడూ~$0$ను ఇచ్చే స్థిర
ప్రమేయం), $g(x_0,y,z)=\Add(z,x_0)$ (తన చివరి, మొదటి ఆర్గుమెంటుల
మొత్తాన్ని ఇచ్చే $3$-స్థానిక ప్రమేయం):
\begin{align*}
  \Mult(x_0,0) & =  f(x_0) = 0 \\
  \Mult(x_0,y+1) & =  g(x_0, y, \Mult(x_0,y)) =
  \Add(\Mult(x_0,y), x_0)
\end{align*}

\end{document}
